% Solution: Tadelis (2013) Exercise 10.9 "Negative Externalities". Input by \solitem.
% New for the Lecture 5 list. Check: check_L05.py.
\begin{exsol}{Tadelis Exercise 10.9}{218--219}{Foundations}{2}
\asked Two firms impose external costs on each other: find the stage-game Nash equilibrium and
compare it with the Pareto-optimal outputs; then find the discount factors for which the Pareto-optimal
outputs are an SPE of the infinitely repeated game.

\method Grim trigger and the one-stage deviation principle with normalised payoffs (Lecture~5,
slides~21, 24--25, 28--29; Tadelis \S10.3, pp.~196--202, threshold (10.5), p.~202). The
cooperative state is checked against the best one-period deviation followed by the stage Nash
equilibrium forever; the punishment state needs no check because it is a stage equilibrium.

\solution
Payoffs: $v_1 = (30 - q_2)q_1 - q_1^2$ and $v_2 = (30 - q_1)q_2 - q_2^2$.

\textbf{(a) Stage game.} The first-order condition of firm~1 is $30 - q_2 - 2q_1 = 0$, so the best
responses are
\[
  q_1(q_2) = \tfrac{30 - q_2}{2}, \qquad q_2(q_1) = \tfrac{30 - q_1}{2}
\]
(downward-sloping lines, each with intercept $15$ on its own axis and $30$ on the rival's axis, crossing at one point). Solving, $q_1 = q_2 = 10$ is the unique Nash equilibrium, with
profits $v_i = (30-10)\cdot10 - 100 = 100$.

\emph{Pareto optimum.} Let $s = q_1 + q_2$. Then
$v_1 + v_2 = 30s - 2q_1q_2 - q_1^2 - q_2^2 = 30s - s^2$, which depends on $s$ only and is maximal at
$s = 15$. Any split of $15$ maximises joint profit; the symmetric one is $q_1 = q_2 = \frac{15}{2}$, with
\[
  v_i = \bigl(30 - 7.5\bigr)\,7.5 - 7.5^2 = 168.75 - 56.25 = 112.5 .
\]
Compared with Nash, each firm produces less ($7.5 < 10$) and earns more ($112.5 > 100$). Each firm
ignores the harm its output does to its neighbour, so the Nash equilibrium overproduces.

\textbf{(b) Grim trigger.} Strategy: produce $7.5$ as long as both have; otherwise produce $10$ forever.
\begin{itemize}
  \item \emph{Punishment state.} The continuation is the stage Nash equilibrium whatever happens, and a
  deviation from $10$ cannot raise the stage payoff, so no deviation pays.
  \item \emph{Cooperative state.} The best deviation against $q_j = 7.5$ is the best response
  $q_i = \frac{30-7.5}{2} = 11.25$, with profit $v^d = 11.25\cdot(22.5 - 11.25) = 11.25^2 = 126.5625$.
\end{itemize}
In normalised payoffs the incentive constraint is
\[
  112.5 \;\ge\; (1-\delta)\,126.5625 + \delta\cdot 100
  \iff \delta \;\ge\; \frac{126.5625 - 112.5}{126.5625 - 100} = \frac{14.0625}{26.5625} = \frac{9}{17}
  \approx 0.529 .
\]
(This is the general threshold (10.5), $\delta^* = (v^d - v^c)/(v^d - v^*)$.) By symmetry the same holds
for firm~2. So the Pareto-optimal outputs $(7.5, 7.5)$ are an SPE outcome of the infinitely repeated
game iff $\delta \ge \frac{9}{17}$.

\interp The externality makes cooperation valuable ($112.5$ against $100$) but the deviation
temptation is large too ($126.5625$). Firms must be fairly patient, more than $0.53$.

\verify \texttt{check\_L05.py} solves the best responses with sympy ($q^* = (10,10)$, $v = 100$,
$\partial^2 v_i/\partial q_i^2 = -2$), shows $v_1 + v_2 = 30s - s^2$, gets $112.5$, $q_d = 45/4$,
$v^d = 2025/16$ and $\delta^* = 9/17$, and tests the inequality on both sides of $9/17$.

\answer{(a) Nash equilibrium $q_1 = q_2 = 10$, profits $100$ each; the Pareto-optimal (joint-profit-maximising)
symmetric outputs are $7.5$ each with profits $112.5$: lower output, higher profit.
(b) Grim trigger sustains $(7.5,7.5)$ iff $\delta \ge \frac{9}{17} \approx 0.53$ (best deviation $q = 11.25$,
profit $126.5625$).}
\end{exsol}
